\subsection{Search strategy}
To answer the research question a systematic literature review was conducted following the PRISMA guidelines \autocite{prisma}. The PRISMA guidelines ensure reproducibility as well as transparency for the selection process of relevant literature. \par 
Search strings were applied to Google Scholar as followed: \textit{''micronutrient§'' AND ''ambient storage'' AND ''fortified'' -fruits} as well as another search string \textit{"micronutrient§ stability" OR "micronutrient§ retention" AND "food fortification"}.
Also, search results were filtered for publications published between 2020 and 2026. Only peer-reviewed literature was considered. The detailed steps are shown in the PRISMA-Flowchart.



\subsection{Inclusion and exclusion criteria}
Studies were included if they were relevant to the research question directly or at least added value to the overall topic of micronutrient retention. Only publications written in English or German were considered. To limit the research scope regarding the broad variety of different foods, it was decided to exclude the term ''fruits'', since works with this keyword adress mainly the pre-fortification stage. As stated in the theoretical framework this is not an exclusion criterion per se but combined with their short shelf-life, fruits were considered unsuitable to answer the research question in its core principle. Consequently, including them would weaken the comparability of findings regarding the post-fortification stage.
%since the aim of this review was to get a deeper understanding of fortified baking goods under different storage conditions. Hence, acceptable food vehicles of fortification are goods like flour, rice, millet and similar food vehicles, since they are on the one hand easy to undergo fortification and on the other hand have a higher shelf life under ambient storage, considering microbiological growth.






\subsection{Data Management}
Since this paper is written by multiple authors, a Citavi-Cloud was created for literature review. Information exchange was facilitated by the Cloud-Service Sciebo and Google-Sheets.
